\documentclass[10pt]{article}
\usepackage[margin=0.65in]{geometry}
\usepackage{booktabs}
\usepackage[T1]{fontenc}
\setcounter{table}{1}
\begin{document}
\begin{table}[ht]
\centering
\caption{Local C2P-Cache workload corpus: 16 complete-compatible traces plus eight later V100-trace extensions.}
\label{tab:c2p-local-workloads}
\small
\begin{tabular}{@{}lll|lll@{}}
\toprule
\textbf{Benchmark} & \textbf{Suite} & \textbf{Abbr.} & \textbf{Benchmark} & \textbf{Suite} & \textbf{Abbr.} \\
\midrule
BFS & ISPASS & BV & sgemm & Parboil & SG \\
LIB & ISPASS & LI & stencil & Parboil & ST \\
LPS & ISPASS & LP & 2DConvolution & PolyBench & 2D \\
RAY & ISPASS & RA & 3mm & PolyBench & 3M \\
b+tree & Rodinia 3.1 & B+ & atax & PolyBench & AT \\
dwt2d & Rodinia 3.1 & DW & bicg & PolyBench & BI \\
gaussian & Rodinia 3.1 & GA & gemm & PolyBench & GE \\
hotspot1 & Rodinia 3.1 & HO & gesummv & PolyBench & GS \\
lud & Rodinia 3.1 & LU & color\_max & Pannotia & CO \\
nn & Rodinia 3.1 & NN & fw\_block & Pannotia & FW \\
cutcp & Parboil & CU & mis & Pannotia & MI \\
mri & Parboil & MR & pagerank & Pannotia & PA \\
\bottomrule
\end{tabular}
\par\medskip
\footnotesize The table preserves the paper's 24-workload naming/abbreviation scheme while recording the local trace corpus.
\end{table}
\end{document}
